\section{Conclusion}

ReversaFeynman extends reverse documentation engineering into an evidence-governed software engineering architecture for agentic maintenance. The original Reversa contribution remains the foundation for extracting traceable operational specifications from legacy systems. ReversaFeynman adds explicit epistemic states, falsifiability-oriented gates, governed human teach-back, adaptive policy evaluation, Hermes-based memory and evidence collection, deterministic evidence governance, and a Software Engineering Intelligence layer for structural localization, isolated execution, multi-candidate repair, quality gates, observability, benchmarking, durable workflows, and shadow-first optimization.

The architectural thesis is that adaptive systems should be allowed to become better at choosing actions without being allowed to redefine what counts as direct evidence. In practical terms, memory, confidence, reward, benchmark performance, static-analysis scores, mutation scores, and optimizer rankings can influence routing and prioritization, but they remain distinct from traceable execution evidence.

The implementation is accompanied by versioned specifications, regression tests, strict contracts, and a proposed experimental protocol. The next research step is empirical rather than architectural: populate ReversaBench with reproducible legacy-specification and issue-repair tasks, run controlled ablations against the original Reversa and minimal repair baselines, and quantify whether the additional layers improve specification reliability, repair success, calibration, regression avoidance, and cost effectiveness.

Accordingly, the current contribution should be read as an implemented research framework and a falsifiable evaluation plan. A claim of superiority is intentionally deferred until the proposed experiments produce sufficient evidence.